\bigskip

\textbf{Pages 12, 13, 14 --- Proximal Policy Optimization (PPO)}

All three pages are about one algorithm: PPO. Think of PPO as a smarter, safer way to train an AI agent.

\subsection*{Page 12 --- PPO's First Two Tricks}

Before the tricks, here is the starting point (the surrogate objective):

\[
\tilde{J}(\theta') \approx \sum_{t,i} \frac{\pi_{\theta'}(a_{t,i} \mid s_{t,i})}{\pi_\theta(a_{t,i} \mid s_{t,i})} \hat{A}^{\pi_\theta}(s_{t,i}, a_{t,i})
\]

\textbf{What each symbol means:}
\begin{itemize}
    \item $\tilde{J}(\theta')$ --- the score we are trying to maximize for the new policy
    \item $\theta'$ --- the parameters (settings) of the new policy we are training
    \item $\theta$ --- the parameters of the old policy we collected data with
    \item $\pi_{\theta'}(a_{t,i} \mid s_{t,i})$ --- probability the new policy picks action $a$ in situation $s$
    \item $\pi_\theta(a_{t,i} \mid s_{t,i})$ --- probability the old policy picks the same action
    \item $\hat{A}^{\pi_\theta}(s_{t,i}, a_{t,i})$ --- the advantage: how much better this action was compared to average
    \item $\sum_{t,i}$ --- add this up across all time steps $t$ and all collected episodes $i$
\end{itemize}

\textbf{What this equation is saying:}
We collected experience using the old policy $\theta$, but we want to improve our new policy $\theta'$ without collecting new data each time. The fraction $\frac{\pi_{\theta'}}{\pi_\theta}$ is called the importance weight --- it corrects for the fact that data was collected under a different policy. If the new policy would have chosen the same action even more often, this ratio is big; if less often, it is small. We multiply this ratio by the advantage $\hat{A}$, which tells us whether that action was actually good or bad. If an action was good (positive advantage) AND the new policy would do it more (ratio $> 1$), the score goes up. If the new policy has drifted far from the old one, the ratio explodes to huge or tiny numbers --- that is the problem this equation alone cannot fix, which is why we need the tricks.

\subsubsection*{Trick \#1 --- Clip the importance weights}

\[
\tilde{J}(\theta') \approx \sum_{t,i} \mathrm{clip}\!\left(\frac{\pi_{\theta'}(a_{t,i} \mid s_{t,i})}{\pi_\theta(a_{t,i} \mid s_{t,i})},\ 1-\epsilon,\ 1+\epsilon\right) \hat{A}^{\pi_\theta}(s_{t,i}, a_{t,i})
\]

\textbf{New symbols:}
\begin{itemize}
    \item $\mathrm{clip}(\cdot,\ 1-\epsilon,\ 1+\epsilon)$ --- a safety clamp: if the value goes below $1-\epsilon$, force it to $1-\epsilon$; if it exceeds $1+\epsilon$, force it to $1+\epsilon$
    \item $\epsilon$ (epsilon) --- a small number (typically 0.1 or 0.2) that sets how far the ratio is allowed to stray from 1
\end{itemize}

\textbf{What this equation is saying:}
The importance weight ratio can get dangerously large or small when the new and old policies differ a lot, making training unstable. Clipping simply forces that ratio to stay inside the window $[1-\epsilon,\ 1+\epsilon]$ --- for example $[0.8,\ 1.2]$ if $\epsilon = 0.2$. If the new policy has already changed too much from the old one, we stop rewarding it for changing further. It is like saying ``you have wandered far enough from where you started --- no extra credit for going further.'' This keeps training stable.

\subsubsection*{Trick \#2 --- Take the minimum}

\[
\tilde{J}(\theta') \approx \sum_{t,i} \min\!\left(\frac{\pi_{\theta'}}{\pi_\theta}\hat{A},\ \mathrm{clip}\!\left(\frac{\pi_{\theta'}}{\pi_\theta},\ 1-\epsilon,\ 1+\epsilon\right)\hat{A}\right)
\]

\textbf{New symbols:}
\begin{itemize}
    \item $\min(\cdot, \cdot)$ --- take whichever of the two values is smaller
\end{itemize}

\textbf{What this equation is saying:}
Trick \#1 clipped the ratio, but that alone has a loophole: when the advantage is negative (meaning the action was bad), clipping might accidentally still give the policy an incentive to make changes. Taking the minimum of the clipped and unclipped objective plugs that loophole. Imagine two numbers: the ``raw'' score and the ``clamped'' score. We always take the worse one. This makes PPO pessimistic by design --- it never rewards the policy for making a change that is not clearly safe and beneficial. This is the core PPO objective.

\subsection*{Page 13 --- Trick \#3: Generalized Advantage Estimation (GAE)}

\subsubsection*{N-step Advantage}

\[
\hat{A}^\pi_n(s_t, a_t) = \sum_{t'=t}^{t+n} \gamma^{t'-t}\, r(s_{t'}, a_{t'}) - \hat{V}^\pi_\phi(s_t) + \gamma^n \hat{V}^\pi_\phi(s_{t+n})
\]

\textbf{What each symbol means:}
\begin{itemize}
    \item $\hat{A}^\pi_n(s_t, a_t)$ --- the estimated advantage of taking action $a_t$ at state $s_t$, looking $n$ steps ahead
    \item $\sum_{t'=t}^{t+n}$ --- sum rewards from now ($t$) up to $n$ steps in the future
    \item $\gamma^{t'-t}$ --- discount factor: rewards further in the future count less (closer = worth more)
    \item $r(s_{t'}, a_{t'})$ --- actual reward received at each step
    \item $\hat{V}^\pi_\phi(s_t)$ --- the critic's estimated value of the current state (what it expected)
    \item $\gamma^n \hat{V}^\pi_\phi(s_{t+n})$ --- the discounted value estimate at the state $n$ steps later (a bootstrapped guess of what comes after)
\end{itemize}

\textbf{What this equation is saying:}
Instead of looking either one step ahead (fast but inaccurate) or all the way to the end (accurate but noisy), this looks exactly $n$ steps ahead. You sum up the real rewards you collected for $n$ steps, then add the critic's guess of what happens after step $n$ (instead of running all the way to the end). Subtracting $\hat{V}(s_t)$ turns this from a raw return into an advantage --- how much better or worse things went compared to what the critic expected. Small $n$ = lower variance but possibly biased; large $n$ = more accurate but noisier. GAE blends all values of $n$ together.

\subsubsection*{GAE (Generalized Advantage Estimation)}

\[
\hat{A}^\pi_\mathrm{GAE}(s_t, a_t) = \sum_{n=1}^{\infty} w_n\, \hat{A}^\pi_n(s_t, a_t), \quad w_n \propto \lambda^{n-1}
\]

\textbf{What each symbol means:}
\begin{itemize}
    \item $\hat{A}^\pi_\mathrm{GAE}$ --- the final blended advantage estimate
    \item $\sum_{n=1}^{\infty}$ --- add up the advantage for every possible horizon: 1-step, 2-step, 3-step\ldots all the way out
    \item $w_n$ --- the weight (importance) given to the $n$-step advantage
    \item $\lambda$ (lambda) --- a number between 0 and 1 that controls how fast the weights shrink
    \item $\lambda^{n-1}$ --- means 1-step gets weight 1, 2-step gets weight $\lambda$, 3-step gets $\lambda^2$, etc.
\end{itemize}

\textbf{What this equation is saying:}
Instead of committing to one fixed horizon $n$, GAE takes a weighted average of all horizons at once. The 1-step advantage gets the biggest weight, the 2-step gets slightly less, the 3-step even less, and so on --- shrinking by a factor of $\lambda$ each time. When $\lambda = 0$, only 1-step matters (low variance, more bias). When $\lambda = 1$, all horizons matter equally (low bias, high variance). You tune $\lambda$ to get the best of both worlds. In practice $N \ll T$ (you do not actually go to infinity --- you cut off early), and cutting earlier reduces variance at the cost of a little bias.

\subsection*{Page 14 --- The Full PPO Algorithm}

\begin{enumerate}
    \item \textbf{Collect data:} Run the current policy $\pi_\theta$ in the environment to gather many episodes.
    \item \textbf{Fit critic:} Train the value network $\hat{V}^\pi_\phi$ to match the observed summed rewards.
    \item \textbf{Compute GAE advantage:} Use the $n$-step formula with weights $w_n \propto \lambda^{n-1}$ to get $\hat{A}^\pi_\mathrm{GAE}$ for every action taken.
    \item \textbf{Update policy:} Run $M$ gradient steps (not just one!) on the full clipped and min surrogate objective from page 13.
\end{enumerate}

The key payoff of all this machinery is Step 4: because PPO has safe clipping, you can squeeze out multiple gradient updates from a single batch of data without the policy flying off the rails. That is something vanilla policy gradients could not do --- they had to collect new data after every single update.

\subsection*{The Big Picture Across All Three Pages}

Think of it like a student who took one test (collected data under old policy $\theta$). A naive student studies that test to death and overcorrects --- completely changing how they think. PPO is the careful student: they use clipping to say ``I will only shift my thinking so far from what I believed before the test,'' use GAE to get smart credit assignment for each answer, and then review the same test paper multiple times safely --- getting more learning out of each exam without going off the rails.

\bigskip

\textbf{Give me a concrete example of trick 3 to explain it.}

\subsection*{Concrete Example: GAE (Trick \#3)}

\textbf{Scenario:} A robot is learning to walk. At moment $t$, it steps forward. We want to know: was that step a good action?

\textbf{Given numbers} ($\gamma = 1$, no discounting, to keep it clean):

\begin{center}
\begin{tabular}{ll}
\hline
 & Value \\
\hline
Reward at step $t$ & $+2$ (moved forward) \\
Reward at step $t+1$ & $+1$ (another step) \\
Reward at step $t+2$ & $-8$ (fell down) \\
Critic's guess for $V(s_t)$ & $3$ \\
Critic's guess for $V(s_{t+1})$ & $4$ \\
Critic's guess for $V(s_{t+2})$ & $5$ \\
Critic's guess for $V(s_{t+3})$ & $0$ \\
\hline
\end{tabular}
\end{center}

\subsubsection*{Step 1 --- Compute each $n$-step advantage}

\textbf{1-step} (looks only 1 step ahead):
\[
A_1 = r_t + V(s_{t+1}) - V(s_t) = 2 + 4 - 3 = +3
\]
``Stepped forward, got $+2$, next state looks good. Action seems great!''

\textbf{2-step} (looks 2 steps ahead):
\[
A_2 = r_t + r_{t+1} + V(s_{t+2}) - V(s_t) = 2 + 1 + 5 - 3 = +5
\]
``Two good steps, state ahead looks even better. Looks amazing!'' --- but the fall is not visible yet.

\textbf{3-step} (looks 3 steps ahead):
\[
A_3 = r_t + r_{t+1} + r_{t+2} + V(s_{t+3}) - V(s_t) = 2 + 1 + (-8) + 0 - 3 = -8
\]
``Now we can see the robot falls at step $t+2$. The action looks terrible.''

\subsubsection*{Step 2 --- GAE blends all three with $\lambda$-weights}

Weights: $w_1 = 1,\quad w_2 = \lambda,\quad w_3 = \lambda^2$

\[
A_\mathrm{GAE} = 1 \cdot A_1 + \lambda \cdot A_2 + \lambda^2 \cdot A_3
\]

\begin{center}
\begin{tabular}{llll}
\hline
$\lambda$ value & Meaning & Calculation & Result \\
\hline
$\lambda = 0$ & Only 1-step & $1(+3)$ & $+3$ --- blind to the fall \\
$\lambda = 0.5$ & Balanced & $1(+3) + 0.5(+5) + 0.25(-8)$ & $+3.5$ --- fall partially visible \\
$\lambda = 0.9$ & Looks far ahead & $1(+3) + 0.9(+5) + 0.81(-8)$ & $+1.02$ --- fall heavily weighted \\
$\lambda = 1$ & All equal weight & $1(+3) + 1(+5) + 1(-8)$ & $0$ --- fall fully cancels the gains \\
\hline
\end{tabular}
\end{center}

\subsubsection*{The punchline}

The 1-step alone would tell the robot ``stepping forward was a great move ($+3$)!'' --- because it cannot see the fall that happens two steps later.

The 3-step alone is too noisy --- it blames the entire step forward harshly ($-8$) even though the first two steps were fine.

GAE with $\lambda = 0.9$ says $+1.02$ --- which is honest: the step forward was slightly good, but not as good as it looked short-term, because trouble is coming. That is exactly the kind of precise, blame-accurate signal the actor needs to learn well.